% !TeX root = main.tex

\renewcommand{\abstractname}{摘要}
\begin{abstract}
在非侵入式词级语言解码中，联合建模邻近词窗可提高预测性能，但依赖目标词之后的输入也会增加输出等待。
本文面向已知词级对齐的自然连续语音聆听任务，提出一种有限输入等待下的局部上下文与滚动修正方法，
使同一词位能够先输出初判，再随最多两个后续词窗到达逐步更新，并固定更早的预测。
实验采用法语脑磁（Pallier2025）和中文脑电（ChineseEEG-2）数据。
在固定 $50$ 词候选集上，半径为 $2$ 的局部上下文相较完整上下文模型，
宏平均 Top-$10$ 分别提高 $5.10$ 和 $5.12$ 个百分点；
共享滚动模型从 K0 更新至 K2 时，该指标分别提高 $19.06$ 和 $14.53$ 个百分点。
序列解码以法语 MEG 为主要分析场景：滚动策略使平均片段 BLEU-1 从 $8.37\%$ 提高至 $13.46\%$，
而固定初始表示的文本重搜仅产生较小改善。
中文 EEG 上的文本收益主要体现于 WER 和词位准确率。
滚动策略相对首次预测平均增加约半秒的神经输入等待，比完整片段策略平均减少约 $1.5$--$1.8$ 秒。
上述结果表明，有限回修可在既定词级对齐与序列搜索协议下提供质量—等待折中。
\end{abstract}
